\section{Introduction}
\label{sec:intro}

Cloud Identity and Access Management (IAM) systems are the core control plane for modern cloud infrastructure. In environments such as Amazon Web Services (AWS), security boundaries are governed by declarative JSON policies that specify fine-grained authorization rules across entities (users, roles, and groups), actions, and resource identifiers~\cite{aws2023iam}. However, managing permissions across tens of thousands of dynamic cloud workloads is notoriously difficult. Misconfigurations---ranging from overly permissive administrative wildcards to subtle unauthorized service linkages---frequently result in devastating privilege escalation, compliance violations, and catastrophic data breaches.

Traditional defenses rely heavily on static rule-based analysis (e.g., Cloud Custodian~\cite{cloudcustodian2023} or AWS IAM Access Analyzer) or formal methods (e.g., SMT-based policy checkers like Zelkova~\cite{backes2018semantic}). While static linters effectively catch obvious violations such as bare \texttt{"Action": "*"} statements, they suffer from two major operational limitations. First, they are rigid and brittle: writing comprehensive rules to capture context-dependent excessive privileges across thousands of distinct AWS API actions is labor-intensive and error-prone. Second, rule engines are blind to peer context; they cannot discern whether a moderately broad permission pattern is customary for a specialized data engineering role or an anomalous privilege anomaly in a web application tier.

To overcome these hurdles, recent security research proposed graph representation learning for unsupervised anomaly detection in cloud IAM. By modelling the interaction topology of cloud entities, attached policies, permitted actions, and target resources as a heterogeneous graph, structural node embeddings (specifically Node2Vec~\cite{grover2016node2vec}) can map access control configurations into a low-dimensional metric space. An unsupervised density estimator, such as Local Outlier Factor (LOF)~\cite{breunig2000lof}, can then score policies as anomalous when their embedding diverges from the cluster formed by standard organizational configurations.

Despite the conceptual promise of graph-embedding anomaly detection, evaluating and deploying these techniques in production faces severe practical obstacles. Prior academic proposals evaluated their methods on private, proprietary corporate snapshots that cannot be shared or peer-reviewed. Furthermore, critical experimental details---such as out-of-sample embedding projection, data contamination handling, hyperparameter sensitivity, and head-to-head comparisons against non-graph baselines---remain largely unexamined. Consequently, cloud practitioners and site reliability engineers lack answers to crucial questions: \emph{Does Node2Vec genuinely capture structural permission semantics beyond basic action counts? How resilient is anomaly detection when training data contains existing misconfigurations? And which types of IAM privilege escalation can actually be detected?}

\subsection{Contributions}
In this paper, we conduct an independent re-implementation, systematic reproduction, and empirical evaluation of graph-embedding anomaly detection for AWS IAM misconfigurations. Our key contributions are:

\begin{enumerate}
    \item \textbf{Open Synthetic IAM Benchmark}: We release a fully reproducible, deterministic benchmark generator synthesizing 3 distinct enterprise organization profiles (\texttt{alpha}, \texttt{beta}, \texttt{gamma}) populated with realistic role hierarchies and 6 traceable misconfiguration mutation operators (P1--P3 privilege escalations and A1--A3 excessive actions), complemented by realistic hard negatives.
    \item \textbf{Leak-Free Graph Architecture}: We build a faithful, leak-free pipeline implementing Entity--Policy--Action--Resource graph construction, seeded out-of-sample Node2Vec embeddings, four anomaly detectors (LOF, One-Class SVM, Isolation Forest, Robust Covariance), and validation-calibrated decision thresholds.
    \item \textbf{Comprehensive Baseline Comparison and Ablation}: We rigorously compare Node2Vec against prevalence, wildcard, permissiveness heuristics, hand-crafted structural features, and bag-of-actions baselines across 30 seeded repetitions with bootstrap confidence intervals, isolating the true contribution of graph random walks.
    \item \textbf{In-Depth Failure Analysis}: We present a granular per-operator and per-difficulty error breakdown, documenting where graph embeddings reliably detect anomalies (easy privilege escalation operators like wildcard assignments) and diagnosing where and why they fail (hard contextual mutations and subtle cross-service attachments).
\end{enumerate}

